\section{施加 \texorpdfstring{$N(\A)$}{N(A)}-等变}
\markboth{施加 \texorpdfstring{$N(\A)$}{N(A)}-等变}{}

范畴 $D(\Bun\rt_B)$ 应当是双商
\[
  B(K)\backslash G(\A)/G(\mathbb O)
\]
上函数空间的几何类比。

现在引入一个范畴，记作 $D(\overline{\Bun}_T)$，它是双商
\[
  T(K)N(\A)\backslash G(\A)/G(\mathbb O)
\]
上函数空间的几何类比。

注意，正如在集合论意义上有
\[
  B(K)\backslash G(\A)/G(\mathbb O)\simeq B(K)\backslash B(\A)/B(\mathbb O),
\]
我们也有
\[
  T(K)N(\A)\backslash G(\A)/G(\mathbb O)\simeq T(K)\backslash T(\A)/T(\mathbb O).
\]
然而，正如范畴 $D(\Bun\rt_B)$ 不等价于 $D(\Bun_B)$，$D(\Bun\rt_T)$ 也不会等同于 $D(\Bun_T)$。

\ob{2.1}{blk:2.1}\textbf{范畴。}叠 $D(\bBun_B)$ 带一个典范群胚的作用，记作 $\mathcal H_N$，它实现 $N(\A)$ 的 adelic 作用。读者可参阅 [Ga1] 或 [Ga2]，那里对此有详细解释。\footnote{原文把「范畴」写成了「叠」：带作用的是范畴 $D(\bBun_B)$，$\mathcal H_N$ 是作用在其上的群胚。译文照录原文，读者按范畴理解即可。}

令 $D(\bBun_T)$ 为 $D(\bBun_B)$ 中 $\mathcal H_N$-等变对象的范畴。由于 $\mathcal H_N$ 投射幂幺，遗忘函子
\[
  D(\bBun_T)\to D(\bBun_B)
\]
全忠实。我们把这个函子记作 $\bfQ^*$。它有右伴随，记作 $\bfQ_*$。

\emph{注。}尽管除 $GL_2$ 的情形外我们不能把 $\bBun_T$ 定义为叠，为记号方便我们仍如此处理。

\ob{2.1.1}{blk:2.1.1}容易看出群胚 $\mathcal H_{G,B}$ 与 $\mathcal H_N$ 的作用交换。因此，考虑 $D(\Bun\rt_B)$ 中 $\mathcal H_N$-等变对象的范畴是合理的；它是 $D(\Bun\rt_B)$ 的全子范畴，我们记作 $D(\Bun\rt_T)$。

等价地，注意作用在 $D(\bBun_B)$ 上的单子 $\Av_{\mathcal H_{G,B}}$ 把 $D(\bBun_T)$ 送到自身。我们可以把 $D(\Bun\rt_T)$ 描述为 $D(\bBun_T)$ 中 $\Av_{\mathcal H_{G,B}}$-模的范畴。

把相应的全忠实嵌入记作
\[
  (\fQ\rt)^* : D(\Bun\rt_T)\hookrightarrow D(\Bun\rt_B),
\]
其右伴随记作 $(\fQ\rt)_*$。

记
\[
  \ind_T : D(\bBun_T)\rightleftarrows D(\Bun\rt_T) : \oblv_T
\]
为所得的伴随函子对。

\ob{2.2}{blk:2.2}\textbf{与 $D(\Bun_T)$ 的关系。}开嵌入
\[
  \fJ_B : \Bun_B\hookrightarrow\bBun_B
\]
诱导开嵌入
\[
  \fJ_T : \Bun_T\hookrightarrow\bBun_T,
\]
以及相应的函子
\[
  (\fJ_T)^! : D(\bBun_T)\to D(\Bun_T).
\]

\begin{othm}{引理}{2.2.1}{thm:2.2.1}
函子 $(\fJ_T)^!$ 有左伴随。
\end{othm}

换言之，引理说的是：$(\fJ_B)^!$ 的（部分定义的）左伴随在全子范畴 $D(\Bun_T)\subset D(\Bun_B)$ 上有定义。理由如下：对象 $(\fJ_B)_!(k_{\Bun_B})\in D(\bBun_B)$（因完整（holonomic）而有定义）关于投影 $\bfQ:\bBun_B\to\bBun_T$ 是 ULA 的。\footnote{这里 $k_Y$ 记概形或叠 $Y$ 上的「常值层」D-模。}

把 $(\fJ_T)^!$ 的左伴随记作 $(\fJ_T)_!$。

\ob{2.2.2}{blk:2.2.2}考虑复合函子
\[
  (\fJ\rt_T)_! := \ind_T\circ(\fJ_T)_! : D(\Bun_T)\to D(\Bun\rt_T),
\]
及其右伴随
\[
  (\fJ\rt_T)^! := \fJ_T^!\circ \oblv_T .
\]
容易看出函子 $(\fJ\rt_T)^!$ 保守。因此范畴 $D(\Bun\rt_T)$ 可以描述为 $D(\Bun_T)$ 上关于单子
\[
  \widetilde\Omega := (\fJ\rt_T)^!\circ(\fJ\rt_T)_!
  \simeq (\fJ_T)^!\circ\oblv_T\circ\ind_T\circ(\fJ_T)_!
\]
的模。（记号 $\widetilde\Omega$ 源自 [GL:Eis] 的记号约定。）

\begin{annote}
$(\fJ\rt_T)^!$ 保守（即它只把同构映到同构）是「单子模范畴」描述的关键：正因为保守，$D(\Bun\rt_T)$ 才能由它在 $D(\Bun_T)$ 上的像与单子 $\widetilde\Omega$ 完全刻画。证明的思路是：$(\fJ\rt_T)^!$ 是 $\bfQ^*$ 与 $\fJ_T^!$ 的复合，前者在等变范畴上是「取底层对象」，后者在 $\Bun_T$ 上是限制，二者在相应范畴上都保守。\par\vspace{0.15em}{\raggedleft 【来源：原书 p.4 §2.2.2（A）；译者按（上下文推断，无外部依据）】\par}
\end{annote}

\ob{2.3}{blk:2.3}\textbf{单子的描述。}单子 $\widetilde\Omega$ 可以显式描述。对元素 $\lambda\in\Lambda$，令
\[
  \fJ_B^\lambda : X^\lambda\times\Bun_B\to\bBun_B
\]
为相应分层的局部闭嵌入。我们用 $\fJ_T^\lambda$ 记号地表示「局部闭嵌入」$X^\lambda\times\Bun_T\to\bBun_T$；我们的意思只是：存在函子
\[
  (\fJ_T^\lambda)^! : D(\bBun_T)\to D(X^\lambda\times\Bun_T).
\]
如同引理 2.2.1，函子 $(\fJ_T^\lambda)^!$ 有左伴随，我们记作 $(\fJ_T^\lambda)_!$。

考虑函子
\[
  \widetilde\Omega^\lambda : D(\Bun_T)\to D(\Bun_T),
\]
它由
\[
  \bigl(H(X^\lambda,-)\boxtimes \mathrm{Id}_{D(\Bun_T)}\bigr)\circ(\fJ_T^\lambda)^!\circ(\fJ_T)_!
\]
给出。

由构造，函子
\[
  \bigoplus_{\lambda\in\Lambda^{pos}}\widetilde\Omega^\lambda
\]
带有自然的单子结构。

\begin{othm}{引理}{2.3.1}{thm:2.3.1}
上述单子 $\bigoplus_{\lambda\in\Lambda^{pos}}\widetilde\Omega^\lambda$ 典范同构于 $\widetilde\Omega$。
\end{othm}

引理来自同构
\[
  \Bun_B\underset{\bBun_B}{\times}\mathcal H_{G,B}\simeq\bigsqcup_\lambda X^\lambda\times\Bun_B,
\]
其中「第二投影」$\Bun_B\times_{\bBun_B}\mathcal H_{G,B}\to\bBun_B$ 由映射 $\bigsqcup_\lambda\fJ_B^\lambda$ 给出。

\begin{annote}
为何只对 $\lambda\in\Lambda^{pos}$ 求和：$\mathcal H_{G,B}$ 的两点只在一般点处相等，从第一个约化变到第二个约化所需的除子必为非负组合，故参数 $\lambda$ 落在正半群 $\Lambda^{pos}$ 中。$X^\lambda$ 是曲线 $X$ 的相应部分对称化幂；当 $\lambda=0$ 时 $X^0=\mathrm{pt}$，$\widetilde\Omega^0$ 就是恒等函子。\par\vspace{0.15em}{\raggedleft 【来源：原书 §2.3 与 [GL:contr] 中 $\Bun_B\times_{\bBun_B}\mathcal H_{G,B}$ 的分解（A）】\par}
\end{annote}
